\scanpage{44}\mspage{6}

Donc il reste essentiellement \deleted{\unsure{3|8}}6 quantités, i.e. \deleted{pour} il suffit, pour \deleted{\unsure{le|tt}} \deleted{\unsure{partitio|partition}} \deleted{subdivision $(2,2)$ \unsure{ti} de $I$ ; dis. $I'\sqcup I''$,} partie à deux éléments $\underset{\substack{\Vert\\\{i'_1,i'_2\}}}{I'}\subset \underset{\substack{\Vert\\\{i''_1,i''_2\}}}{I^{\unsure{2|*|\prime\prime}}}$, de de déterminer l'une des quatre quantités

\[
 \rho_{\underset{\tikz{\draw[<->] (0,0) to[bend right=35] (.36,0);}}{i'_1\,i'_2}\;
 \underset{\tikz{\draw[<->] (0,0) to[bend right=35] (.36,0);}}{i''_1\,i''_2}}.
 \qquad\deleted{\unsure{Il}}
\]

$\star$ Nous allons \deleted{\unsure{don}} enfin donner une formule pour relier

\[
\rho_{ij\unsure{k|h}\unsure{l|\ell}}\quad\text{et}\quad\rho_{i\deleted{\unsure{j}}kjl}
\]

(ce qui permettra donc de passer d'un repère à un repère quelconque)

\[
 \rho_{ijkl}+\rho_{ikjl}=1
\]

i.e

\[
\boxed{\rho_{i\deleted{\unsure{jk}}kjl}=1-\rho_{i\deleted{\unsure{kj}}jkl}}\tag*{(7)}
\]

\margin{[ Faire vérification}

\uline{Prop} Soit un système de $u_{ijkl}:V_i\longrightarrow V_j$, satisfaisant les relations (2) (3) ; d'où (4) des $\rho_{ijkl}\in\operatorname{Aut}(V_i)$, \unsure{)}, satisfaisant (5)(6). Pour que ça corresponde à une \unsure{quadrité|quadrite|quadralité} \deleted{\unsure{$V$}} $(V,\deleted{\unsure{\varphi}}(V_i),(\varphi_i))$ \deleted{\unsure{$\subset$}} $V\subset\prod V_i$, il

% archive sheet number: 43
